\scanpage{5}
7) Modules pour $S\subset K\subset X$

On suppose \inserted{\deleted{$X$} compact, et} pour simplifier $X$ \deleted{compact} \inserted{sans bord} orienté.
Ceci dit, on fixe $S_0\subset K_0\subset X_0$ situation
top, et on \uline{rigidifie} $(S\subset K\subset X)$ en
se donnant une classe d'isotopie d'homeo.
top. avec la situation standart. On élimine
ainsi les autos, sauf dans le cas
\deleted{$\unsure{K_0}$} où $X_0$ a des composantes de degré
$0,1$\ldots

Pour avoir un ``\unsure{nb} fini de modules'',
on \inserted{pourrait} exige\inserted{r} que les arêtes de $K$ soient
des arcs de cercle pour la conn.
projective canonique. Mais les
``cartes \inserted{coll.} holomorphes'' ne satisfont pas
néc. à cette condition ! On pourrait
exiger plutôt que pour \deleted{\unsure{nos}} les
paramétrisations géodésiques sur le
découpé, \unsure{ses composés|ses diagrammes} avec l'homéomorphisme
de recollage soient linéaires. [Éventuellement
on pourrait exiger l'équidistance \inserted{des sommets} \deleted{\unsure{sur}}
chaque composante connexe du bord de
le découpé\ldots] Mais on a des ennuis
pour les morceaux \deleted{de genre $0$} \inserted{qui sont des disques \unsure{non ponctués}}
: pas de métrique nat \unsure{const.}. Il
% archive foliation: 4
